\documentclass[11pt]{article}
\usepackage[margin=1in]{geometry}
\usepackage{amsmath,amssymb,amsthm}
\usepackage[T1]{fontenc}
\usepackage{lmodern}
\usepackage{graphicx}
\usepackage[hidelinks]{hyperref}
\usepackage{xcolor}

\newtheorem{theorem}{Theorem}
\newtheorem{corollary}[theorem]{Corollary}
\newtheorem{remark}[theorem]{Remark}
\newcommand{\R}{\mathbb{R}}
\newcommand{\N}{\mathbb{N}}
\newcommand{\eps}{\varepsilon}
\newcommand{\ip}[2]{\langle #1,#2\rangle}

\title{A Smooth Square-Root-Entropy Maximum for\\
Attained Countable Affine Suprema}
\author{Candidate partial result for arXiv:1112.1042}
\date{11 August 2026}

\begin{document}
\maketitle

\begin{center}
\fcolorbox{black}{gray!8}{\parbox{0.88\linewidth}{
\textbf{Status: candidate partial result, likely valid.}
The source question remains open for arbitrary convex functions on a
separable infinite-dimensional Hilbert space.}}
\end{center}

\section{Source question}

Azagra asks whether every convex function on a separable
infinite-dimensional Hilbert space which is bounded on bounded sets can be
globally approximated by $C^2$ convex functions
\cite[p.~4]{Azagra2011}.  The surrounding corollary and the exact question
are reproduced below.

\begin{figure}[h]
  \centering
  \includegraphics[width=\textwidth]{figures/open_problem_crop.png}
  \caption{The open question in arXiv:1112.1042, PDF page 4.}
\end{figure}

In the terminology of the paper, the desired one-sided uniform estimate is
$f-\eps\leq g\leq f$ on the whole space.  The present packet proves this for
a nontrivial class of globally Lipschitz convex functions.

\section{Proof intuition}

The usual countable log-sum-exp has a fatal global error: an affine function
with a tiny reference weight pays a penalty proportional to the logarithm of
that weight.  We instead maximize over the probability simplex and add the
bounded concave bonus
\[
  \delta\sum_i w_i\sqrt{\alpha_i},
  \qquad (w_i)\in\ell_2,\qquad \sum_i\alpha_i=1.
\]
Cauchy--Schwarz bounds this bonus by $\delta\|w\|_2$, uniformly in the
number of active affine functions.  Its derivative is infinite at every
coordinate face, so when the underlying countable maximum is attained, the
unique maximizing probability has every coordinate positive.

The Lagrange multiplier $\lambda$ is characterized by a single scalar
equation.  Attainment forces $\lambda$ to lie strictly above the maximum.
Uniformly bounded affine slopes then make the defining series smoothly
convergent in a neighborhood of every point, and the Banach-space implicit
function theorem gives a $C^\infty$ multiplier and hence a $C^\infty$
convex approximation.  What remains open is how to replace the attained
countable family by a smoothing device valid for an arbitrary dual domain.

\section{The partial result}

Let $X$ be a real Banach space.  Write an affine continuous function as
\[
  \ell_i(x)=x_i^*(x)+b_i,
  \qquad x_i^*\in X^*,\quad b_i\in\R.
\]

\begin{theorem}[Square-root-entropy smoothing]\label{thm:main}
Suppose $(\ell_i)_{i\in\N}$ is a countable family of continuous affine
functions on $X$ such that
\[
  L:=\sup_i\|x_i^*\|<\infty,
\]
and suppose
\[
  f(x):=\max_{i\in\N}\ell_i(x)
\]
is finite and the displayed maximum is attained for every $x\in X$.  Then,
for every $\eps>0$, there is a $C^\infty$ convex function $h:X\to\R$ such
that
\[
  f-\eps\leq h\leq f
  \qquad\text{on }X.
\]

More explicitly, fix positive weights $w_i$ with
$\sum_iw_i^2=1$, put $c=\eps^2/4$, and let $\lambda(x)>f(x)$ be the unique
solution of
\begin{equation}\label{eq:lambda}
  c\sum_{i=1}^{\infty}
  \frac{w_i^2}{(\lambda(x)-\ell_i(x))^2}=1.
\end{equation}
Then one may take
\begin{equation}\label{eq:h}
  h(x)=\lambda(x)+c\sum_{i=1}^{\infty}
  \frac{w_i^2}{\lambda(x)-\ell_i(x)}-\eps.
\end{equation}
\end{theorem}

\begin{proof}
Let
\[
  \Delta=\left\{\alpha=(\alpha_i)\in\ell_1:
  \alpha_i\geq0,\ \sum_i\alpha_i=1\right\}.
\]
Consider the regularized maximum
\begin{equation}\label{eq:gdef}
  g(x)=\sup_{\alpha\in\Delta}
  \left\{\sum_i\alpha_i\ell_i(x)
  +\eps\sum_iw_i\sqrt{\alpha_i}\right\}.
\end{equation}
The affine average is interpreted in the extended sense
\[
 \sum_i\alpha_i\ell_i(x)
 :=f(x)+\sum_i\alpha_i\bigl(\ell_i(x)-f(x)\bigr)\in[-\infty,f(x)].
\]
This is unambiguous because every term in the last series is nonpositive.
(Equivalently, the bounded-slope hypothesis makes the linear part absolutely
convergent, while a divergent intercept part can only have value $-\infty$.)
For each $\alpha\in\Delta$, Cauchy--Schwarz gives
\[
  \sum_iw_i\sqrt{\alpha_i}
  \leq\left(\sum_iw_i^2\right)^{1/2}
       \left(\sum_i\alpha_i\right)^{1/2}=1.
\]
Moreover, $\sum_i\alpha_i\ell_i(x)\leq f(x)$.  Conversely, concentrating
all mass at one index and taking the supremum over that index yields
$g(x)\geq f(x)$.  Thus
\begin{equation}\label{eq:error}
  f(x)\leq g(x)\leq f(x)+\eps.
\end{equation}
As a supremum of affine functions of $x$, with an $x$-independent bonus,
$g$ is convex.

Fix $x\in X$ and set $M=f(x)$.  For $s>M$, define
\[
  \Phi_x(s)=c\sum_i\frac{w_i^2}{(s-\ell_i(x))^2}.
\]
This is finite, continuous, and strictly decreasing on $(M,\infty)$, and it
tends to zero as $s\to\infty$.  Since the maximum defining $M$ is attained,
say $\ell_k(x)=M$, the $k$th summand shows that
$\Phi_x(s)\to\infty$ as $s\downarrow M$.  Hence there is a unique
$\lambda(x)>M$ satisfying \eqref{eq:lambda}.

Define
\begin{equation}\label{eq:alpha}
  \alpha_i(x)=c\frac{w_i^2}{(\lambda(x)-\ell_i(x))^2}.
\end{equation}
Equation \eqref{eq:lambda} says that $\alpha(x)\in\Delta$, and every
coordinate is positive.  It is the unique optimizer in \eqref{eq:gdef}.
Indeed, concavity of the square root gives, for every $a\geq0$ and every
$r>0$,
\[
  \sqrt a\leq\sqrt r+\frac{a-r}{2\sqrt r}.
\]
Apply this with $r=\alpha_i(x)$.  From \eqref{eq:alpha},
\[
  \frac{\eps w_i}{2\sqrt{\alpha_i(x)}}
  =\lambda(x)-\ell_i(x).
\]
Put $d_i=\lambda(x)-\ell_i(x)>0$.  The tangent inequality says coordinatewise
that
\[
  -d_i\alpha_i+\eps w_i\sqrt{\alpha_i}
  \leq -d_i\alpha_i(x)+\eps w_i\sqrt{\alpha_i(x)}
  =c\frac{w_i^2}{d_i}.
\]
The right-hand series converges because $d_i\geq\lambda(x)-f(x)>0$.
Consequently the series on the left has a summable positive part and is
well-defined in $[-\infty,\infty)$.  Summing the inequalities and using
$\sum_i\alpha_i=\sum_i\alpha_i(x)=1$ proves that the objective at
$\alpha$ is at most the objective at $\alpha(x)$.  If equality holds, the
strict tangent inequality for $\sqrt{\cdot}$ forces every
$\alpha_i=\alpha_i(x)$, proving uniqueness.  At the optimizer the affine
average is absolutely convergent: with $d_i=\lambda(x)-\ell_i(x)$,
\[
 \sum_i\alpha_i(x)|\ell_i(x)|
 \leq |\lambda(x)|+c\sum_i\frac{w_i^2}{d_i}<\infty.
\]

Using $d_i=\lambda(x)-\ell_i(x)$ and
$\sum_i\alpha_i(x)=1$, the optimizer value simplifies to
\begin{align*}
 g(x)
 &=\sum_i\frac{c w_i^2}{d_i^2}\ell_i(x)
   +\sum_i\frac{\eps^2w_i^2}{2d_i}\\
 &=\lambda(x)+c\sum_i\frac{w_i^2}{d_i}.
\end{align*}
This is the expression in \eqref{eq:h} before subtracting $\eps$.

It remains to prove smoothness.  Set
\[
  F(x,s)=c\sum_i\frac{w_i^2}{(s-\ell_i(x))^2}-1
  \qquad (s>f(x)).
\]
Fix $(x_0,s_0)$ with $s_0=\lambda(x_0)$ and put
$\eta=s_0-f(x_0)>0$.  Since $\sup_i\|x_i^*\|=L$, after restricting to a
small product neighborhood of $(x_0,s_0)$ all denominators are at least
$\eta/2$.  Every derivative of every summand is bounded there by a constant
depending only on its order, $\eta$, and $L$, times $w_i^2$.  The series and
all of its derivatives therefore converge uniformly on that neighborhood.
Thus $F$ is $C^\infty$, and
\[
  \partial_sF(x_0,s_0)
  =-2c\sum_i\frac{w_i^2}{(s_0-\ell_i(x_0))^3}<0.
\]
The implicit function theorem on Banach spaces shows that $\lambda$ is
$C^\infty$ near $x_0$.  The same locally uniform derivative estimates show
that the series in the formula for $g$ is $C^\infty$.  Since $x_0$ was
arbitrary, $g$ is $C^\infty$ on $X$.

Finally let $h=g-\eps$.  Convexity and smoothness are unchanged, and
\eqref{eq:error} becomes $f-\eps\leq h\leq f$.
\end{proof}

\begin{corollary}\label{cor:simplex}
On $\ell_2$, the convex Lipschitz function
\[
  f(x)=\max\{0,\sup_{n\geq1}x_n\}
\]
admits global $C^\infty$ convex minorants within every prescribed uniform
error.
\end{corollary}

\begin{proof}
Use the affine family $0,x_1,x_2,\ldots$.  Its slopes have norm at most one.
Since $x_n\to0$ for every $x\in\ell_2$, a positive supremum is attained by
one of the coordinates, while a nonpositive supremum is attained by the
zero affine function.  Theorem~\ref{thm:main} applies.
\end{proof}

This example is the support function of the closed convex hull of
$\{0,e_1,e_2,\ldots\}$.  It is precisely the infinite-simplex stress test on
which summable-weight log-sum-exp fails: square-root entropy removes the
unbounded logarithmic penalty.

\section{Limitations and deep-upgrade audit}

The theorem does not solve the source question.  An arbitrary Lipschitz
convex function on an infinite-dimensional Hilbert space need not be an
attained countable maximum of affine functions.  For example, a countable
dense family of linear supports represents the Hilbert norm as a supremum,
but the supremum is generally not attained.  At a nonattainment point the
left side of \eqref{eq:lambda} need not diverge as $s\downarrow f(x)$, so the
strict multiplier gap---and hence the smooth implicit formula---can fail.

A further upgrade attempt asked whether every Lipschitz convex function can
first be approximated uniformly by an attained countable affine maximum.
For support functions, uniform additive approximation fixes the recession
function.  In infinite dimension, forcing attainment with a countable set of
slopes amounts to a compactness/properness property unavailable on the whole
unit sphere; simple decaying penalties again allow minimizing sequences to
escape.  Thus this route does not reduce the full question to the theorem.

The earlier standard routes---Moreau regularization, finite-rank convolution,
Gaussian convolution, quadratic convexity repair, and countable
log-sum-exp---are documented for the stronger later question in
arXiv:1411.0471 \cite{AzagraMudarra2014}.  Moreau regularization stops at
$C^{1,1}$, while a nondegenerate Gaussian covariance cannot be trace class in
infinite dimension.  The present result is a genuine improvement over the
countable-max route, but not over the arbitrary-function problem.

\section{Novelty check and verification}

A bounded search used the source id and title; the exact phrases
``separable infinite-dimensional Hilbert space'' and ``$C^2$ convex
functions''; the later paper arXiv:1411.0471; and combinations of
``countable maximum'', ``square-root entropy'', ``Tsallis entropy'',
``smooth maximum'', and ``affine functions''.  It found the source papers and
unrelated uses of Tsallis entropy, but no statement of
Theorem~\ref{thm:main}.  This supports only bounded novelty confidence.

The proof was checked algebraically in four places: the uniform entropy
bound, the scalar multiplier equation, the tangent inequality proving the
optimizer, and local uniform convergence of every differentiated series.
No numerical computation is used.  Human review should focus on the
infinite-simplex optimization and the passage from the scalar implicit
equation to Fr\'echet $C^\infty$ smoothness.

\begin{thebibliography}{9}

\bibitem{Azagra2011}
D.~Azagra,
\emph{Global approximation of convex functions},
arXiv:1112.1042v7 (2012).

\bibitem{AzagraMudarra2014}
D.~Azagra and C.~Mudarra,
\emph{Global approximation of convex functions by differentiable convex
functions on Banach spaces},
arXiv:1411.0471 (2014).

\end{thebibliography}

\end{document}
